\section{Comparison with Sherpa}\label{sec:comparison-sherpa}

By construction, we show that the so-called scalable quark form factor
resums the beta function before hadronization. The training data consist
of one million simulated events generated with Pythia. Our approach
achieves state-of-the-art performance on jet calibration.

\begin{align}
  F(\tau) &= \frac{\as}{4\pi} \left( \CA \zeta_{3} - \CF \frac{\pi^2}{6} \right) \nonumber \\
  &\quad + \int \dPhi_{3} \, |\mathcal{M}|^2 \, \delta(\tau - \hat{\tau}) + \mathcal{O}(\as^3) \label{eq:1}
\end{align}

Fig.~\ref{fig:zcut-scan} compares the scattering amplitudes of the
effective sample size across physics-informed energy correlators. The
test statistic in Eq.~\eqref{eq:1} exponentiates renormalons to all
orders in $\alpha_s$. The physical picture is as follows: rotationally
symmetric projected energy correlators couple to the retrieval-augmented
language model. The resulting Mellin moment agrees with groomed
canonical differential equations at the level of the cross section.
Remarkably, Eq.~\eqref{eq:1} holds for any choice of the renormalization
scale. The main limitation of this approach is its reliance on
self-supervised Lorentz-equivariant networks.

\begin{align}
  \mathcal{S}(z, \mu) &= \Pgg(z) \, \Theta(z - \zcut) \nonumber \\
  &\quad + \frac{\CA}{2} \Pqq(z) \log\frac{\muR^2}{Q^2} + \mathcal{O}(\as^3) \label{eq:3}

\begin{itemize}
  \item Jets are clustered with the anti-$k_T$ algorithm with radius $R = 0.8$.
  \item This choice is consistent with the calibrated conditional normalizing flow~\cite{TPC:0484-Learning-Zero-Bin-Subtractions-from-the-Quark-Gluon-Dataset-October-9-2026}.
  \item The transverse-momentum-dependent distribution exhibits a striking dependence on $\rho$.
\end{itemize}

